\documentclass[11pt]{article}
\usepackage[margin=1in]{geometry}
\usepackage{amsmath,amssymb,amsthm}
\usepackage{graphicx}
\usepackage{microtype}
\usepackage[colorlinks=true,linkcolor=blue,citecolor=blue,urlcolor=blue]{hyperref}

\newtheorem{theorem}{Theorem}
\newtheorem{lemma}{Lemma}
\newtheorem{corollary}{Corollary}
\newcommand{\K}{\mathbb K}
\newcommand{\I}{\mathcal I}
\newcommand{\Sp}[1]{\widehat{\otimes}^{#1,s}_{\pi_s}}
\newcommand{\Tp}{\widehat{\otimes}_{\pi}}

\title{Symmetric tensor stability implies tensor stability}
\author{Candidate full solution for arXiv:1603.00246}
\date{11 August 2026}

\begin{document}
\maketitle

\begin{center}
\fbox{\parbox{0.91\linewidth}{\textbf{Status.} Candidate full proof, likely
valid, requiring human review. The answer to the source question is
affirmative. In fact, stability under symmetric squares alone implies full
tensor stability.}}
\end{center}

\section{Source question}

Let $0<p\le1$ and let $(\I,\|\cdot\|_{\I})$ be a $p$-normed operator ideal
over $\K=\mathbb R$ or $\mathbb C$. Botelho and Torres call $\I$
\emph{symmetrically tensorstable} (or \emph{s-tensorstable}) if, for every
$m\in\mathbb N$, there is $C_m>0$ such that
\[
 \otimes^{m,s}u\in
 \I\bigl(\Sp{m}E;\Sp{m}F\bigr),\qquad
 \|\otimes^{m,s}u\|_{\I}\le C_m\|u\|_{\I}^{m}
\]
for all Banach spaces $E,F$ and $u\in\I(E;F)$.

They call $\I$ \emph{tensorstable} if there are constants $D_m>0$ such that
for every $u_j\in\I(E_j;F_j)$,
\[
 u_1\otimes\cdots\otimes u_m\in
 \I(E_1\Tp\cdots\Tp E_m;F_1\Tp\cdots\Tp F_m)
\]
and
\[
 \|u_1\otimes\cdots\otimes u_m\|_{\I}
 \le D_m\prod_{j=1}^m\|u_j\|_{\I}.
\]
After proving the converse implication, they ask at the end of the paper:
\begin{quote}
``We do not know if every s-tensorstable operator ideal is tensorstable.''
\end{quote}
This is the open question on PDF page 27 (printed page 26) of
\cite{BT2018}.

\begin{figure}[ht]
\centering
\includegraphics[width=\linewidth]{figures/open_problem_crop.png}
\caption{The complete open question in \cite[p.~27]{BT2018}, rendered and
cropped from the local source-paper PDF.}
\end{figure}

\section{Proof intuition}

Put two ideal operators on the diagonal of one operator between finite
$\ell_1$-sums. The symmetric square of this diagonal operator has three
coordinate types: the square of the first block, the square of the second
block, and the mixed block. That mixed block is exactly the ordinary
projective tensor product of the two original operators.

There are uniformly bounded maps which insert an ordinary tensor into this
mixed symmetric block and extract it again. Thus symmetric-square stability
of the diagonal operator forces the ordinary tensor product to belong to the
ideal. Normalizing the two diagonal blocks before adding them converts the
$p$-triangle estimate into the required product bound. Once the
two-operator estimate is available, associativity gives arbitrary finite
tensor products by iteration.

\section{The mixed component of a symmetric square}

For a Banach space $X$, write
\[
 z\vee z'=\frac12(z\otimes z'+z'\otimes z)\in\Sp{2}X.
\]
This is the symmetrization convention used in \cite[Proposition~7.12]{BT2018}.

\begin{lemma}\label{lem:mixed}
Let $E,G,F,H$ be Banach spaces. Put $X=E\oplus_1G$ and $Y=F\oplus_1H$,
with canonical coordinate injections $i_E,i_G,i_F,i_H$ and projections
$q_E,q_G,q_F,q_H$. There are bounded operators
\[
 J:E\Tp G\longrightarrow\Sp{2}X,
 \qquad
 Q:\Sp{2}Y\longrightarrow F\Tp H
\]
such that
\[
 J(x\otimes g)=i_E(x)\vee i_G(g),\qquad
 Q(i_F(f)\vee i_H(h))=f\otimes h,
\]
and $\|J\|\le2$, $\|Q\|\le2$.
\end{lemma}

\begin{proof}
Let $S_X^2:X\times X\to\Sp{2}X$ be the symmetrization map
$S_X^2(z,z')=z\vee z'$. Its linearization on $X\Tp X$ has norm equal to
the second polarization constant, which is at most $2$; this is precisely
the estimate used in \cite[Proposition~7.12]{BT2018}. Therefore
\[
 J=(S_X^2)_L\circ(i_E\otimes i_G)
\]
has norm at most $2$ and has the required action on elementary tensors.

For completeness, the same bound follows directly without identifying any
two tensor norms. The polarization identity
\[
 z\vee z'=\frac14\bigl((z+z')^{\otimes2}-(z-z')^{\otimes2}\bigr)
\]
and replacement of $(z,z')$ by $(tz,t^{-1}z')$, followed by minimization over
$t>0$, give
\[
             \pi_s(z\vee z')\le2\|z\|\,\|z'\|.
\]
The universal property of $\Tp$ then yields the asserted extension of $J$.

Let
\[
 \iota_Y^2:\Sp{2}Y\longrightarrow Y\Tp Y
\]
be the canonical norm-one inclusion used in
\cite[Proposition~7.12]{BT2018}, and define
\[
                 Q=2(q_F\otimes q_H)\circ\iota_Y^2.
\]
Then $\|Q\|\le2$. Moreover,
\[
 \iota_Y^2(i_F(f)\vee i_H(h))
 =\frac12\bigl(i_F(f)\otimes i_H(h)+i_H(h)\otimes i_F(f)\bigr).
\]
The coordinate tensor projection $q_F\otimes q_H$ keeps the first term and
annihilates the second, so the factor $2$ gives
$Q(i_F(f)\vee i_H(h))=f\otimes h$.
\end{proof}

\section{Affirmative solution}

\begin{theorem}\label{thm:main}
Let $0<p\le1$. Every s-tensorstable $p$-normed operator ideal is
tensorstable. More precisely, suppose only that symmetric squares satisfy
\[
 \otimes^{2,s}w\in\I(\Sp{2}X;\Sp{2}Y),\qquad
 \|\otimes^{2,s}w\|_{\I}\le C_2\|w\|_{\I}^{2}
\]
for all $w\in\I(X;Y)$. Then, for all $u\in\I(E;F)$ and
$v\in\I(G;H)$,
\begin{equation}\label{eq:pair-bound}
 \|u\otimes v\|_{\I}
 \le K\|u\|_{\I}\|v\|_{\I},
 \qquad K=2^{,2+2/p}C_2.
\end{equation}
Consequently one may take $D_m=K^{m-1}$ in the definition of
tensorstability.
\end{theorem}

\begin{proof}
The conclusion is immediate if $u=0$ or $v=0$, so suppose both are nonzero
and set
\[
                         a=\|u\|_{\I},\qquad b=\|v\|_{\I}.
\]
Use the spaces and coordinate maps of Lemma~\ref{lem:mixed} and define the
block-diagonal operator
\[
 W=i_F\frac{u}{a}q_E+i_H\frac{v}{b}q_G:X\longrightarrow Y.
\]
The ideal property puts both summands in $\I(X;Y)$ with ideal $p$-norm at
most one. Since $\I$ is $p$-normed,
\begin{equation}\label{eq:W-bound}
                    \|W\|_{\I}^{p}\le1+1=2,
 \qquad\text{hence}\qquad \|W\|_{\I}^{2}\le2^{2/p}.
\end{equation}

The defining formula for $\otimes^{2,s}W$ on pure squares, polarized, gives
\begin{equation}\label{eq:polarized-action}
 (\otimes^{2,s}W)(z\vee z')=Wz\vee Wz'
 \qquad(z,z'\in X).
\end{equation}
Thus, for $x\in E$ and $g\in G$, Lemma~\ref{lem:mixed} and the diagonal
form of $W$ yield
\begin{align*}
 Q(\otimes^{2,s}W)J(x\otimes g)
 &=Q\left(i_F\frac{u(x)}a\vee i_H\frac{v(g)}b\right)\\
 &=\frac{u(x)}a\otimes\frac{v(g)}b.
\end{align*}
By density of elementary tensors,
\begin{equation}\label{eq:factorization}
 Q\circ(\otimes^{2,s}W)\circ J
 =\frac{u}{a}\otimes\frac{v}{b}.
\end{equation}

Symmetric-square stability, the operator-ideal property,
Lemma~\ref{lem:mixed}, and \eqref{eq:W-bound} now give
\begin{align*}
 \left\|\frac{u}{a}\otimes\frac{v}{b}\right\|_{\I}
 &\le\|Q\|\,\|\otimes^{2,s}W\|_{\I}\,\|J\|\\
 &\le 2\,C_2\|W\|_{\I}^{2}\,2\\
 &\le 2^{,2+2/p}C_2=K.
\end{align*}
Multiplication by $ab$ proves \eqref{eq:pair-bound}.

Finally apply \eqref{eq:pair-bound} repeatedly. Identifying iterated
projective tensor products through their canonical isometric associativity
maps, induction gives
\[
 \|u_1\otimes\cdots\otimes u_m\|_{\I}
 \le K^{m-1}\prod_{j=1}^m\|u_j\|_{\I},
\]
and membership in $\I$ at each step. Hence $\I$ is tensorstable.
\end{proof}

\begin{corollary}
For a $p$-normed operator ideal, s-tensorstability and tensorstability are
equivalent. Consequently condition (c) of
\cite[Theorem~7.8]{BT2018} may equivalently be replaced by tensorstability.
\end{corollary}

\begin{proof}
Theorem~\ref{thm:main} proves the previously open implication. The converse
is \cite[Proposition~7.12]{BT2018}.
\end{proof}

\section{Verification, scope, and novelty}

The proof is noncomputational. Its load-bearing points are the $p$-triangle
estimate \eqref{eq:W-bound}, the two uniform mixed-component maps of
Lemma~\ref{lem:mixed}, and the exact intertwining identity
\eqref{eq:factorization}. The accompanying \texttt{verification.md} checks
these steps separately, including conventions, scalar fields, completions,
and iteration. The constant is intentionally safe and is not asserted to be
optimal.

\paragraph{Bounded novelty check.}
On 11 August 2026, the local run indexes and the downloaded arXiv source
corpus were searched by arXiv id, title, authors, and the exact phrases
``s-tensorstable'', ``symmetrically tensorstable'', and the verbatim open
question. The terminology occurred only in the source paper in the local
corpus. External exact-phrase and title/author searches likewise found the
source but no separate answer. OpenAlex metadata listed nine works citing
the journal article; their titles and available records did not identify an
answer. A directly relevant open 2020 dissertation on associativity of
projective and injective tensor products was downloaded and text-searched and
contained none of the exact terminology. No prior statement of
Theorem~\ref{thm:main} or of this block-diagonal mixed-component proof was
found. Novelty confidence is moderate rather than definitive; specialist
review of non-indexed literature is recommended.

\paragraph{Human-review recommendation.}
Check the averaged convention for $\vee$ and hence the factor $2$ in $Q$;
check that the source's norm-one map $\iota_Y^2$ has the stated direction;
and check whether later operator-ideal literature records this equivalence
under different terminology. No unresolved mathematical dependency is known.

\begin{thebibliography}{9}
\bibitem{BT2018}
G. Botelho and E. R. Torres,
\emph{Two-sided polynomial ideals on Banach spaces},
J. Math. Anal. Appl. \textbf{462} (2018), 900--914,
\href{https://doi.org/10.1016/j.jmaa.2017.12.054}{doi:10.1016/j.jmaa.2017.12.054};
circulated as \emph{Polynomial ideals from a nonlinear viewpoint},
\href{https://arxiv.org/abs/1603.00246}{arXiv:1603.00246}.
\end{thebibliography}

\end{document}
